\section{Empirical Strategy}\label{sec:strategy}

The strategy exploits province-level meat-price variation as an
external salience treatment.
The treatment is measured by the National Bureau of Statistics and
merged from outside the CFPS survey, so the household's current
expectation never appears on both sides of any regression.
I estimate three linked equations that together identify the
overreaction pattern.

\subsection{Treatment: Province-Level Meat-Price Shocks}
\label{sec:meatshock}

The treatment variable is the cumulative log change in provincial
meat CPI between two consecutive CFPS waves:
\begin{equation}\label{eq:meat_shock}
  \mathrm{MeatShock}_{p,\,t-1\to t}
  = \sum_{m \in \mathcal{M}(t-1,t)}
      \log\!\left(\frac{\mathrm{MeatCPI}_{pm}}{100}\right),
\end{equation}
where $\mathrm{MeatCPI}_{pm}$ is the NBS sub-index for meat and
poultry in province $p$ and month $m$, and $\mathcal{M}(t-1,t)$
denotes the set of months between the two consecutive survey waves.
The index is published monthly at the province level starting
January 2016, which covers four CFPS wave pairs (2016--2018,
2018--2020, 2020--2022, and the backward link from 2014--2016).

Three properties of this measure matter for identification.
First, it is external to the survey: it contains no information
derived from respondents themselves.
Second, it is cumulative: a household surveyed in 2020 experienced
the full path of meat-price movements since the 2018 wave, not just
the end-point level.
Third, it varies at the province level, which allows me to absorb
region-by-wave fixed effects and still retain within-region
identifying variation.

The 2018--2019 African Swine Fever episode provides the sharpest
source of variation.
ASF is a contagious hemorrhagic disease in pigs with near-100
percent mortality in infected herds; it has no transmission pathway
to humans and no direct biological channel to non-pork markets.
Provincial herd losses from ASF varied substantially because culling
policy, transport restrictions, and initial herd density differed
across provinces.
This heterogeneity generates cross-sectional dispersion in the
treatment that is plausibly orthogonal to province-level demand
conditions, given the zoonotic nature of the shock.

\subsection{Three-Equation System}
\label{sec:three_equations}

\paragraph{Equation 1: expectations.}
I estimate how meat-price shocks shift household inflation
expectations:
\begin{equation}\label{eq:eq1}
  \mu_{ipt}
  = \alpha
  + \beta_1\,\mathrm{MeatShock}_{p,\,t-1\to t}
  + \gamma' X_{ipt}
  + \delta_{r(p)\times t}
  + \varepsilon_{ipt},
\end{equation}
where $\mu_{ipt} \in \{-1,0,1\}$ is the ordinal inflation
expectation of household $i$ in province $p$ at wave $t$,
$X_{ipt}$ is a vector of household controls, and
$\delta_{r(p)\times t}$ are region-by-wave fixed effects.
Standard errors are clustered at the province level.
A positive $\hat{\beta}_1$ indicates that households in provinces
with larger meat-price increases report higher inflation expectations.

\paragraph{Equation 2: pass-through.}
I estimate how much of the meat-price shock passes through
to forward headline CPI:
\begin{equation}\label{eq:eq2}
  \pi^{\mathrm{headline}}_{p,\,t\to t+1}
  = \alpha
  + \beta_2\,\mathrm{MeatShock}_{p,\,t-1\to t}
  + \delta_{r(p)\times t}
  + u_{pt},
\end{equation}
where $\pi^{\mathrm{headline}}_{p,\,t\to t+1}$ is realised
province-level headline CPI inflation over the subsequent inter-wave
period.
This regression runs at the province-wave level.
The coefficient $\beta_2$ is the reduced-form pass-through: how much
headline inflation follows from a unit meat-price shock.

\paragraph{Equation 3: forecast errors.}
The reduced-form implication of overreaction is that households in
high-shock provinces subsequently overpredict inflation:
\begin{equation}\label{eq:eq3}
  \mathrm{FE}_{ipt}
  = \alpha
  + \beta_3\,\mathrm{MeatShock}_{p,\,t-1\to t}
  + \gamma' X_{ipt}
  + \delta_{r(p)\times t}
  + \varepsilon_{ipt},
\end{equation}
where $\mathrm{FE}_{ipt} \equiv \pi^{\mathrm{headline}}_{p,\,t\to t+1}
- \tilde{\mu}_{ipt}$ is realised provincial headline inflation over
the next inter-wave period minus the household's directional
expectation scaled to CPI units.
A negative forecast error means that the household expected more
inflation than materialised (overprediction).
Under overreaction, $\beta_3 < 0$: provinces with larger meat-price
shocks generate more negative forecast errors because expectations
overshoot what realised inflation delivers.

\paragraph{The overreaction test.}
The three coefficients are linked by an approximate decomposition:
\begin{equation}\label{eq:decomp_strategy}
  \beta_3 \approx \beta_2 - k\,\beta_1,
\end{equation}
where $k$ converts the ordinal expectation into implied inflation
units.
If $k\,\beta_1 > \beta_2$, the forecast error is negative.
Under rational expectations, $\beta_1$ should equal the
pass-through $\beta_2$ scaled by meat's CPI weight
$\omega_m \approx 0.05$, which is small.
The overreaction prediction is $\hat{\beta}_1 \gg \hat{\beta}_2$:
households move beliefs by more than the shock warrants given its
actual headline weight.
The sign of $\beta_3$ provides a direct test that requires no
assumption on $k$.

\subsection{Identification Assumptions}
\label{sec:identification_assumptions}

The strategy requires two conditions.

\paragraph{Condition 1: relevance.}
The meat-price shock must shift household inflation expectations
($\beta_1 \neq 0$).
This is testable.
Pork is the most widely consumed meat in China and accounts for a
large share of household food expenditure.
Purchases are frequent, so households receive repeated price signals.
The attention literature finds that frequently purchased goods
generate stronger belief updating than infrequently purchased ones
\citep{Cavallo2017, Gorodnichenko2015}.
The ASF episode generated widely covered news about pork prices in
2019 that was difficult to miss even for inattentive consumers.

\paragraph{Condition 2: conditional exogeneity.}
The meat-price shock must be uncorrelated with province-level
expectation determinants other than through salience, conditional
on region-by-wave fixed effects and household controls.
Region-by-wave fixed effects absorb regional business cycles,
regional monetary policy transmission, and nationwide inflation
shocks that hit all provinces within a region simultaneously.
The remaining identifying variation is within-region cross-province
differences in meat-price shocks over time.

The main threat is a province-specific demand shock that
simultaneously raises meat prices and inflation expectations through
a real-income channel.
I address this in two ways.
First, the ASF episode provides a biologically determined source of
supply variation: the disease spread was driven by initial herd
density and animal movement networks, not by province-level
aggregate demand.
Second, the commodity-placebo test in Section~\ref{sec:placebo}
shows that grain shocks of comparable magnitude do not produce the
same expectation response, which is inconsistent with a generic
demand-side confound.

\subsection{Inference with Few Clusters}
\label{sec:inference}

With 31 province clusters, asymptotic cluster-robust standard
errors can be liberal.
I report wild cluster bootstrap $p$-values throughout, using
the \citet{Webb2023} six-point weight distribution with 9,999
replications.
Point estimates and conventional cluster-robust standard errors
are reported alongside bootstrap $p$-values.
All significance statements in the text refer to bootstrap
$p$-values unless otherwise noted.
